% =========================================================================
\section{Singularidades aisladas: polos}\label{sec:polos}
% =========================================================================

Los puntos donde $P$ se anula son aquellos en los que los coeficientes de la
forma normal \eqref{eq:despeje} pueden dejar de ser analíticos. Para
clasificarlos necesitamos medir \emph{cómo} diverge una función. La
herramienta adecuada procede de la variable compleja, y por eso adoptamos
desde ahora el punto de vista de que $x$ puede tomar valores en $\C$.

% -------------------------------------------------------------------------
\subsection{Definición y caracterizaciones}\label{sec:def-polo}
% -------------------------------------------------------------------------

\emph{Idea intuitiva.} Un polo es un punto donde la función se va a
infinito \emph{como una potencia}, igual que $1/x$ o $1/x^{2}$ en $x=0$. El
\emph{orden} mide lo rápido que diverge: $1/x^{2}$ explota más deprisa que
$1/x$.

\begin{definicion}[Polo de orden $p$]\label{def:polo}
Sea $f$ analítica en un entorno reducido
$U^{*}=\{x:0<\abs{x-x_0}<\rho\}$ de $x_0$. Se dice que $f$ tiene un
\emph{polo de orden $p$} en $x_0$, con $p\in\N$, $p\ge1$, si $f$ no admite
extensión analítica a $x_0$ pero
\begin{equation}\label{eq:polo}
  g(x)\defeq(x-x_0)^{p}f(x)\quad\text{es analítica en } x_0
  \quad\text{y}\quad g(x_0)\neq0 .
\end{equation}
Si $f$ coincide en $U^{*}$ con una función analítica en $x_0$, se dice que
$x_0$ es una \emph{singularidad evitable} (o «polo de orden $0$»).
\end{definicion}

Multiplicar por $(x-x_0)^p$ «compensa» exactamente la divergencia.
Despejando, $f=g/(x-x_0)^p$, y como $g(x)\to g(x_0)\eqqcolon c\neq0$,
\begin{equation}\label{eq:asintotico}
  f(x)\sim\frac{c}{(x-x_0)^{p}},\qquad x\to x_0 .
\end{equation}

\begin{proposicion}[Caracterizaciones del orden]\label{prop:orden}
Sea $f$ analítica en un entorno reducido de $x_0$ y no extensible
analíticamente a $x_0$. Para $p\ge1$, son equivalentes:
\begin{enumerate}[label=\textup{(\alph*)}]
  \item $f$ tiene un polo de orden $p$ en $x_0$;
  \item $p$ es el \emph{menor} entero $m\ge0$ tal que $(x-x_0)^{m}f$ es
        analítica en $x_0$;
  \item $f$ admite en un entorno reducido de $x_0$ un desarrollo de
        \emph{Laurent} con un número finito de potencias negativas, la más
        negativa de las cuales es $(x-x_0)^{-p}$:
        \begin{equation}\label{eq:laurent}
          f(x)=\ubr{\frac{g_0}{(x-x_0)^{p}}+\frac{g_1}{(x-x_0)^{p-1}}+\dots
          +\frac{g_{p-1}}{x-x_0}}{\text{parte singular}}
          +\ubr{g_p+g_{p+1}(x-x_0)+\cdots}{\text{parte analítica}}
        \end{equation}
        con $g_0\neq0$.
\end{enumerate}
\end{proposicion}

\begin{proof}
(a)$\Rightarrow$(c). Escribiendo $g=\sum_{j\ge0}g_j(x-x_0)^{j}$ con
$g_0=g(x_0)\neq0$ y dividiendo por $(x-x_0)^{p}$ se obtiene \eqref{eq:laurent}.
(c)$\Rightarrow$(a). Multiplicando \eqref{eq:laurent} por $(x-x_0)^p$ se
obtiene una serie de potencias convergente $g$ con $g(x_0)=g_0\neq0$.

(a)$\Rightarrow$(b). Si $(x-x_0)^{m}f\eqqcolon h$ fuese analítica con
$m<p$, entonces $g=(x-x_0)^{p-m}h$ se anularía en $x_0$, en contra de
$g(x_0)\neq0$. (b)$\Rightarrow$(a). Sea $m$ el mínimo; $m\ge1$ porque $f$
no es extensible. Si $h=(x-x_0)^{m}f$ cumpliera $h(x_0)=0$, se podría
extraer un factor, $h=(x-x_0)h_1$ con $h_1$ analítica, y
$(x-x_0)^{m-1}f=h_1$ sería analítica, contra la minimalidad. Luego
$h(x_0)\neq0$ y $m=p$.
\end{proof}

\begin{observacion}[Por qué «el menor $p$»]\label{obs:menor-p}
Si $(x-x_0)^{2}f$ es analítica, también lo son $(x-x_0)^{3}f$,
$(x-x_0)^{4}f,\dots$, por ser productos de funciones analíticas: multiplicar
de más siempre «funciona». Por eso el orden es el \emph{mínimo} exponente que
lo consigue, y exigir $g(x_0)\neq0$ en \eqref{eq:polo} es una forma
equivalente de imponer esa minimalidad. Un polo de orden $p$ es, en suma,
una serie de potencias a la que se le permiten potencias negativas, pero
sólo hasta $(x-x_0)^{-p}$.
\end{observacion}

% -------------------------------------------------------------------------
\subsection{Ejemplos y clasificación}\label{sec:ejemplos-polos}
% -------------------------------------------------------------------------

\begin{ejemplo}[Polos simples y dobles]\label{ej:polos}
\leavevmode
\begin{enumerate}[label=(\roman*)]
  \item $f(x)=1/x$ en $x_0=0$: $x\,f=1$ es analítica y no nula, luego polo
        de orden $1$ (\emph{polo simple}).
  \item $f(x)=\dfrac{1}{x^{2}(x-1)}$ tiene dos polos de órdenes distintos.
        \begin{itemize}
          \item En $x_0=0$: $x\,f=\dfrac1{x(x-1)}$ sigue divergiendo, pero
                \begin{equation}\label{eq:ej-doble}
                  x^{2}f=\frac{1}{x-1}\ \text{es analítica en }0,\qquad g(0)=-1\neq0,
                \end{equation}
                luego es un polo de orden $2$, con $f\sim-1/x^{2}$ cerca de $0$.
          \item En $x_0=1$: $(x-1)\,f=1/x^{2}$ es analítica en $1$ y vale
                $1\neq0$: polo de orden $1$, con $f\sim1/(x-1)$.
          \item En cualquier otro punto (por ejemplo, $x_0=2$) $f$ es
                analítica: no hay polo.
        \end{itemize}
  \item $f(x)=\dfrac{\sin x}{x^{3}}$ en $x_0=0$. Parece de orden $3$, pero
        \begin{equation}\label{eq:sen-x3}
          \frac{\sin x}{x^{3}}=\frac{x-\frac{x^{3}}{6}+\frac{x^{5}}{120}-\cdots}{x^{3}}
          =\frac1{x^{2}}-\frac16+\frac{x^{2}}{120}-\cdots ,
        \end{equation}
        de modo que $x^{2}f=1-\tfrac{x^{2}}{6}+\cdots$ es analítica con valor
        $1\neq0$, mientras que $x\,f\sim1/x$ todavía diverge: polo de orden
        $2$. El cero del numerador «se come» una potencia.
  \item $f(x)=\dfrac{\sin x}{x}=1-\dfrac{x^{2}}{6}+\cdots$ no tiene un polo:
        definiendo $f(0)=1$ es analítica (singularidad evitable).
\end{enumerate}
\end{ejemplo}

\begin{proposicion}[Regla práctica para cocientes de polinomios]\label{prop:regla}
Sean $N$, $D$ polinomios sin factores comunes salvo, eventualmente, una
potencia de $(x-x_0)$. Si $x_0$ es un cero de multiplicidad $m$ de $D$ y de
multiplicidad $k\ge0$ de $N$, entonces $N/D$ tiene en $x_0$ un polo de orden
$m-k$ si $m>k$, y una singularidad evitable si $m\le k$.
\end{proposicion}

\begin{proof}
Escribiendo $N=(x-x_0)^{k}\tilde N$ y $D=(x-x_0)^{m}\tilde D$ con
$\tilde N(x_0)\tilde D(x_0)\neq0$, se tiene
$(x-x_0)^{m-k}\,N/D=\tilde N/\tilde D$, analítica y no nula en $x_0$
(\cref{prop:algebra}).
\end{proof}

No toda singularidad es un polo. Las funciones $\sqrt{x}$, $\ln x$ y
$\eul^{1/x}$ en $x_0=0$ no lo son, porque ninguna potencia $x^{p}$ las vuelve
analíticas: $x^{p}\sqrt{x}=x^{p+1/2}$ tiene exponente no entero; $x^{p}\ln x$
conserva el logaritmo; y
\[
  \eul^{1/x}=\sum_{m=0}^{\infty}\frac{1}{m!\,x^{m}}
\]
tiene \emph{infinitas} potencias negativas. Son singularidades «peores» que
un polo. La \cref{tab:singularidades} resume la clasificación.

\begin{table}[ht]
\centering
\caption{Clasificación de las singularidades en $x_0$ según el desarrollo de
Laurent $f=\sum_{j}c_j(x-x_0)^{j}$ y el comportamiento de $\abs{f}$.}
\label{tab:singularidades}
\renewcommand{\arraystretch}{1.4}
\small
\begin{tabular}{@{}l l l l@{}}
\toprule
\textcolor{cian}{\bfseries Tipo} & \textcolor{cian}{\bfseries Laurent} &
\textcolor{cian}{\bfseries Comportamiento} & \textcolor{cian}{\bfseries Ejemplo}\\
\midrule
Evitable              & sin potencias negativas            & $f$ acotada               & $\dfrac{\sin x}{x}$\\[4pt]
Polo de orden $p$     & la más negativa es $(x-x_0)^{-p}$  & $\abs{f}\sim\abs{c}\,\abs{x-x_0}^{-p}$ & $\dfrac{1}{x^{2}(x-1)}$\\[4pt]
Esencial              & infinitas potencias negativas      & sin límite en $\C$        & $\eul^{1/x}$\\[4pt]
Punto de ramificación & no existe desarrollo de Laurent    & $f$ multivaluada en $\C$  & $\sqrt{x}$, $\ln x$\\
\bottomrule
\end{tabular}
\end{table}

La \cref{fig:fases} muestra la firma de cada polo en el plano complejo: en
un \emph{retrato de fase} se colorea cada punto $z$ según el argumento
$\arg f(z)$, y alrededor de un polo de orden $p$ el círculo cromático se
recorre exactamente $p$ veces (en sentido opuesto al de un cero de orden $p$),
de acuerdo con el principio del argumento \cite[cap.~4]{AHLFORS}.

% Figura: retratos de fase de 1/z, 1/z^2 y 1/(z^2 (z-1))
% Se usa sombreado plano (una celda = un color) para no interpolar a través
% del salto de +180 a -180 grados del argumento.
\begin{figure}[ht]
\centering
\begin{tikzpicture}
  \begin{groupplot}[group style={group size=3 by 1, horizontal sep=0.55cm},
      oscuro, grid=none, view={0}{90}, width=5.15cm, height=5.15cm,
      axis equal image, colormap name=fase, shader=flat corner,
      point meta min=-180, point meta max=180,
      xmin=-1.5, xmax=1.5, ymin=-1.5, ymax=1.5,
      xtick={-1,0,1}, ytick={-1,0,1},
      samples=\resF, samples y=\resF, domain=-1.5:1.5, y domain=-1.5:1.5,
      tick label style={font=\tiny}, title style={yshift=-2pt},
      xlabel={$\Real z$}, xlabel style={font=\scriptsize, yshift=4pt}]
    % (a) 1/z: arg = -atan2(y,x)
    \nextgroupplot[title={(a) $1/z$\,: polo simple},
        ylabel={$\Imag z$}, ylabel style={font=\scriptsize, yshift=-4pt}]
      \addplot3[surf, point meta={-atan2(y,x)}] (x, y, 0);
      \addplot3[only marks, mark=x, mark size=3pt, line width=1.2pt, blanco]
        coordinates {(0,0,0)};
    % (b) 1/z^2: arg = -2 atan2(y,x), reducido a (-180,180]
    \nextgroupplot[title={(b) $1/z^{2}$\,: polo doble}, yticklabels={}]
      \addplot3[surf, point meta={mod(-2*atan2(y,x)+540,360)-180}] (x, y, 0);
      \addplot3[only marks, mark=x, mark size=3pt, line width=1.2pt, blanco]
        coordinates {(0,0,0)};
    % (c) 1/(z^2 (z-1)): arg = -2 atan2(y,x) - atan2(y,x-1)
    \nextgroupplot[title={(c) $1/\bigl(z^{2}(z-1)\bigr)$}, yticklabels={},
        colorbar right, colorbar style={width=0.22cm, ytick={-180,0,180},
          yticklabels={$-\pi$,$0$,$\pi$}, tick label style={font=\tiny, text=gris},
          axis line style={grismedio}, ylabel={$\arg f$},
          ylabel style={font=\scriptsize, text=texto, yshift=-2pt}}]
      \addplot3[surf, point meta={mod(-2*atan2(y,x)-atan2(y,x-1)+900,360)-180}] (x, y, 0);
      \addplot3[only marks, mark=x, mark size=3pt, line width=1.2pt, blanco]
        coordinates {(0,0,0) (1,0,0)};
  \end{groupplot}
\end{tikzpicture}
\caption{Retratos de fase: cada punto $z$ se colorea según $\arg f(z)$
(barra de la derecha). Alrededor de un polo de orden $p$ (marcado con
$\times$) el círculo cromático se recorre $p$ veces: una vez en (a), dos en
(b). En (c) conviven un polo doble en $z=0$ y uno simple en $z=1$, que son
los del \cref{ej:polos}\,(ii).}
\label{fig:fases}
\end{figure}


% -------------------------------------------------------------------------
\subsection{Relación con la ecuación diferencial}\label{sec:polos-edo}
% -------------------------------------------------------------------------
Escribamos \eqref{eq:edo} en forma normal con
$p\defeq A/P$ y $q\defeq B/P$. Si $P(x_0)\neq0$, $p$ y $q$ son analíticas en
$x_0$ (\cref{prop:algebra}): no hay nada que compensar. Si $P$ tiene un cero
de orden $m$ en $x_0$, por la \cref{prop:regla} los cocientes pueden tener
polos de orden \emph{a lo sumo} $m$ (menos, si $A$ o $B$ se anulan también
en $x_0$ y cancelan factores).

\begin{ejemplo}[Legendre y Bessel]\label{ej:legendre-bessel}
\leavevmode
\begin{enumerate}[label=(\roman*)]
  \item \emph{Legendre}, $P=1-x^{2}$:
        \[
          p=\frac{-2x}{(1-x)(1+x)},\qquad
          q=\frac{\ell(\ell+1)}{(1-x)(1+x)}
        \]
        tienen polos de orden $1$ en $x=\pm1$ (si $\ell(\ell+1)\neq0$).
  \item \emph{Bessel}, $x^{2}y''+x\,y'+(x^{2}-\nu^{2})\,y=0$:
        \[
          p=\frac1x\ \ (\text{orden }1),\qquad
          q=\frac{x^{2}-\nu^{2}}{x^{2}}\ \ (\text{orden }2\text{ si }\nu\neq0)
        \]
        en $x=0$. Si $\nu=0$, $q\equiv1$ es analítica.
\end{enumerate}
\end{ejemplo}

\begin{observacion}[Avance de la Parte II]\label{obs:avance}
El orden de estos polos es lo que distinguirá, en la segunda parte del tema,
los puntos \emph{singulares regulares} ---en los que $p$ tiene a lo sumo un
polo de orden $1$ y $q$ a lo sumo uno de orden $2$, es decir, $(x-x_0)\,p$ y
$(x-x_0)^{2}q$ son analíticas--- de los \emph{irregulares}. En los primeros
se buscan soluciones de la forma $(x-x_0)^{r}\sum a_n(x-x_0)^n$ (método de
Frobenius), con un exponente $r$ fijado por la \emph{ecuación indicial}
\cite{TEMAV}. Tanto $x=\pm1$ en Legendre como $x=0$ en Bessel son singulares
regulares.
\end{observacion}

% -------------------------------------------------------------------------
\subsection{Por qué los polos complejos limitan el radio}\label{sec:radio}
% -------------------------------------------------------------------------
La función $f(x)=1/(1+x^{2})$ no tiene ninguna singularidad real: es
analítica en todo $\R$, acotada y de clase $\Cinf$. Y, sin embargo, su
desarrollo en $x_0=0$,
\begin{equation}\label{eq:lorentz}
  \frac{1}{1+x^{2}}=1-x^{2}+x^{4}-x^{6}+\cdots=\sum_{n=0}^{\infty}(-1)^{n}x^{2n},
\end{equation}
\emph{sólo converge para $\abs{x}<1$}: es una serie geométrica de razón
$-x^{2}$. La razón está en el plano complejo: $1+z^{2}=0$ en $z=\pm\iu$, que
son polos de orden $1$ situados a distancia $1$ del origen. Una serie de
potencias converge en un \emph{disco}, y ese disco no puede contener un
polo; por eso el radio es la distancia al polo complejo más cercano
(\cref{fig:radio}).

\begin{teorema}[Radio de convergencia y singularidades]\label{teo:radio}
Sea $f$ holomorfa en el disco $\disco{z_0}{\rho}=\{z\in\C:\abs{z-z_0}<\rho\}$.
Entonces la serie de Taylor de $f$ en $z_0$ converge a $f$ en todo
$\disco{z_0}{\rho}$. En particular, si $f$ es racional, el radio de
convergencia de su serie de Taylor en $z_0$ es exactamente la distancia de
$z_0$ al polo más próximo.
\end{teorema}

\begin{proof}
La primera afirmación es el teorema de Cauchy--Taylor \cite[§4.3]{AHLFORS}.
Para la segunda, sea $d$ la distancia al polo más cercano $z_1$; la primera
parte da $R\ge d$. Si fuera $R>d$, la suma de la serie sería continua en el
disco cerrado $\overline{\disco{z_0}{d}}$, y por tanto acotada en él, y
coincidiría con $f$ en $\disco{z_0}{d}$; pero $\abs{f(z)}\to\infty$ cuando
$z\to z_1$ dentro de ese disco. Luego $R=d$.
\end{proof}

% Figura: 1/(1+x^2) — polos complejos ±i y radio de convergencia R = 1
\begin{figure}[ht]
\centering
\begin{tikzpicture}
  % ---------------- (a) Plano complejo: retrato de fase y disco ----------------
  \begin{axis}[name=fase, oscuro, grid=none, view={0}{90},
      width=6.3cm, height=6.3cm, axis equal image,
      colormap name=fase, shader=flat corner,
      point meta min=-180, point meta max=180,
      xmin=-2, xmax=2, ymin=-2, ymax=2,
      xtick={-2,-1,0,1,2}, ytick={-2,-1,0,1,2},
      samples=\resF, samples y=\resF, domain=-2:2, y domain=-2:2,
      xlabel={$\Real z$}, ylabel={$\Imag z$},
      title={(a) Plano complejo: $\arg\dfrac{1}{1+z^{2}}$}]
    \addplot3[surf, point meta={-atan2(2*x*y,1+x^2-y^2)}] (x, y, 0);
    % disco de convergencia |z| < 1
    \addplot3[blanco, line width=1.1pt, domain=0:360, samples=120, samples y=0]
      ({cos(x)}, {sin(x)}, 0);
    % intervalo real de convergencia
    \addplot3[blanco, line width=2.6pt, samples y=0] coordinates {(-1,0,0) (1,0,0)};
    \addplot3[only marks, mark=*, mark size=1.8pt, blanco, samples y=0]
      coordinates {(-1,0,0) (1,0,0)};
    % radio
    \addplot3[blanco, line width=0.8pt, densely dashed, samples y=0]
      coordinates {(0,0,0) (0.7071,-0.7071,0)};
    % polos
    \addplot3[only marks, mark=x, mark size=4pt, line width=1.5pt, blanco, samples y=0]
      coordinates {(0,1,0) (0,-1,0)};
    \node[fill=fondo, fill opacity=0.8, text opacity=1, font=\scriptsize,
          inner sep=1.5pt, rounded corners=1pt] at (axis cs:0.42,1.36,0) {$+\iu$};
    \node[fill=fondo, fill opacity=0.8, text opacity=1, font=\scriptsize,
          inner sep=1.5pt, rounded corners=1pt] at (axis cs:0.42,-1.36,0) {$-\iu$};
    \node[fill=fondo, fill opacity=0.8, text opacity=1, font=\scriptsize,
          inner sep=1.5pt, rounded corners=1pt] at (axis cs:1.35,-0.95,0) {$R=1$};
  \end{axis}

  % ---------------- (b) Recta real: sumas parciales ----------------------------
  \begin{axis}[at={($(fase.east)+(1.45cm,0)$)}, anchor=west, oscuro,
      width=7.2cm, height=5.9cm,
      xmin=-1.6, xmax=1.6, ymin=-0.6, ymax=1.6,
      xtick={-1.5,-1,-0.5,0,0.5,1,1.5}, ytick={-0.5,0,0.5,1,1.5},
      xlabel={$x$}, title={(b) Recta real: sumas parciales $S_N$},
      legend style={at={(0.5,0.03)}, anchor=south, legend columns=2},
      samples=241, domain=-1.6:1.6]
    \fill[grisoscuro, opacity=0.55] (axis cs:-1.6,-0.6) rectangle (axis cs:-1,1.6);
    \fill[grisoscuro, opacity=0.55] (axis cs:1,-0.6) rectangle (axis cs:1.6,1.6);
    \addplot[c2, restrict y to domain=-5:5] {(1+x^6)/(1+x^2)};
    \addplot[c3, restrict y to domain=-5:5] {(1-x^12)/(1+x^2)};
    \addplot[c4, restrict y to domain=-5:5] {(1+x^26)/(1+x^2)};
    \addplot[blanco, line width=1.6pt] {1/(1+x^2)};
    \legend{$S_{4}$, $S_{10}$, $S_{24}$, $\frac{1}{1+x^{2}}$}
  \end{axis}
\end{tikzpicture}
\caption{(a) Retrato de fase de $f(z)=1/(1+z^{2})$. Los polos $z=\pm\iu$
(marcados con $\times$) están a distancia $1$ del origen; el mayor disco
centrado en $0$ que no los contiene es $\abs{z}<1$, cuya intersección con la
recta real es el intervalo $(-1,1)$ (trazo grueso). (b) Las sumas parciales
$S_{2N}(x)=\sum_{n=0}^{N}(-1)^n x^{2n}$ de \eqref{eq:lorentz} se acercan a
$f$ en $\abs{x}<1$ y se separan bruscamente de ella en $\abs{x}>1$
(zonas sombreadas), pese a que $f$ es perfectamente regular sobre $\R$.}
\label{fig:radio}
\end{figure}


Esta es la idea que explota el \cref{teo:fuchs}: para conocer dónde
converge la serie de una solución \emph{real} hay que mirar dónde se anula
$P$ en el plano \emph{complejo}.
